%% appendix-classes -- proofs for the tree classes section, including the
%% closure under intersection of leaf cones (prop:ch-meet).

\section{Proofs for Tree Classes}\label{app:tc}\label{app:classes}

This appendix proves the results of \cref{sec:tc} whose proofs were
deferred, in its notation. For a covariant tree $t$, $\nu_0t=t$ and
$\nu_1t=t^\partial$, so $L^0_t$ is the set of words with a prefix labeled
$\Top$ in $t$, and $L^1_t$ the set of words with a prefix labeled $\Bot$. By
\cref{lem:bg-normalization}(3), duality reverses the covariant order, and it
commutes with the constructor, $f(t_i)_i^\partial=f(t_i^\partial)_i$, since it
exchanges labels position by position.

\subsection{Observation Pins}\label{app:tc-pins}

\begin{proof}[Proof of \cref{lem:tc-pins}]
Let $\operatorname{Pref}(w)$ be the set of prefixes of $w$, including $w$, and
let $\operatorname{Off}(w)=\Sigma^*\setminus(\operatorname{Pref}(w)\cup
w\Sigma^*)$ be the set of words that leave the pin path, that is, the words
$\wpre wk\,j\,v$ with $k<\len w$ and $j\ne w[k]$. The positions of
$\emb_b(w,g)$ are the words of $\operatorname{Pref}(w)$, of which the proper
prefixes of $w$ are labeled $f$; the children $\wpre wk\,j$ off the pin path,
labeled $b$; and $w$ followed by the positions of $g$. A word of
$\operatorname{Off}(w)$ passes a child off the pin path, a proper prefix of $w$
passes no leaf, and an extension of $w$ follows $g$. Hence the four trees of
\cref{eq:tc-pins} have the following cones:
\[
\begin{array}{lll}
 \text{tree}&L^0&L^1\\\hline
 g_1=\emb_\Bot(w,\Top)&w\Sigma^*&\operatorname{Off}(w)\\
 g_2=\emb_\Top(w,f(\Top^n))&\Sigma^*\setminus\operatorname{Pref}(w)&\varnothing\\
 g_3=\emb_\Top(w,\Bot)&\operatorname{Off}(w)&w\Sigma^*\\
 g_4=\emb_\Bot(w,f(\Bot^n))&\varnothing&\Sigma^*\setminus\operatorname{Pref}(w)
\end{array}
\]
We use \cref{thm:lf-representation}(1,2): the cones of $t$ are extension
closed and disjoint, and $t\le u$ iff $L^0_t\subseteq L^0_u$ and
$L^1_u\subseteq L^1_t$.
\begin{itemize}[leftmargin=*]
\item $g_1\le t$ iff $w\Sigma^*\subseteq L^0_t$ and
  $L^1_t\subseteq\operatorname{Off}(w)$. If $w\in L^0_t$, then
  $w\Sigma^*\subseteq L^0_t$ by extension closure;
  $L^1_t$ contains no extension of $w$, by disjointness, and no proper prefix
  of $w$, since otherwise it would contain $w$; so $L^1_t\subseteq\operatorname{Off}(w)$.
  Conversely, $g_1\le t$ gives $w\in L^0_t$.
\item $t\le g_2$ iff $L^0_t\subseteq\Sigma^*\setminus\operatorname{Pref}(w)$, since the
  condition on the cones of sign $1$ is vacuous. By extension closure, $L^0_t$
  meets $\operatorname{Pref}(w)$ iff $w\in L^0_t$.
\item $t\le g_3$ iff $L^0_t\subseteq\operatorname{Off}(w)$ and $w\Sigma^*\subseteq L^1_t$;
  this is the first case with the two cones exchanged.
\item $g_4\le t$ iff $L^1_t\subseteq\Sigma^*\setminus\operatorname{Pref}(w)$, iff
  $w\notin L^1_t$, as in the second case.
\end{itemize}
The pin path has $\len w+1$ positions, there are $(n-1)\len w$ children off
the pin path, and the end of the pin path has at most $n$ children. So each tree has at
most $1+n(\len w+1)$ positions.

For the consequence, let $x$ be a visible variable and $\rho$ an assignment.
The cones of $\rho(x)$ are the covariant cones of $t=\nu_0\rho(x)$, since
$\nu_1\rho(x)=t^\partial$ by \cref{lem:bg-normalization}(1). So an observation
with variable $x$ and word $w$ is equivalent to one of the four inequalities
between $t$ and a tree $g_k$ of \cref{eq:tc-pins}. By
\cref{lem:bg-normalization}(1,2), $g_k\le\nu_0\rho(x)$ iff
$\nu_0g_k\le_c\rho(x)$, and $\nu_0\rho(x)\le g_k$ iff $\rho(x)\le_c\nu_0g_k$.
The tree $\nu_0g_k$ has the domain of $g_k$.
\end{proof}

\subsection{Signed Bounds and Medians}\label{app:tc-signed}

\begin{proof}[Proof of \cref{lem:tc-signed}]
(1) By \cref{lem:bg-signed}, $\rho$ solves $\chi$ iff $\rho^\pm$ solves
$\chi^\pm$. For a lower bound, \cref{lem:bg-normalization}(2) gives
$g\le_c\rho(x)$ iff $\nu_+g\le\nu_+\rho(x)$ iff $\nu_-\rho(x)\le\nu_-g$; these
are the two bounds of \cref{eq:tc-signed-bounds} at $\rho^\pm$. Upper bounds
are symmetric.

(2) Since duality is an involution, so is $\iota$. Let $\sigma$ be an assignment
of $V^\pm$. In every row of \cref{tab:bg-signed}, and for the two bounds of
each pair in \cref{eq:tc-signed-bounds}, $\iota\sigma$ satisfies one entry iff
$\sigma$ satisfies the other, because duality reverses the covariant order,
commutes with the constructor, exchanges $\Bot$ and $\Top$, and maps $\nu_+g$
to $\nu_-g$. For instance, $\iota\sigma$ satisfies $u^+\le v^+$ iff
$\sigma(u^-)^\partial\le\sigma(v^-)^\partial$, iff $\sigma(v^-)\le\sigma(u^-)$;
it satisfies $f(a_i^{c(i)})_i\le u^+$ iff
$f(\sigma(a_i^{\bar c(i)})^\partial)_i\le\sigma(u^-)^\partial$, iff
$\sigma(u^-)\le f(\sigma(a_i^{\bar c(i)}))_i$; it satisfies $u^+=\Bot$ iff
$\sigma(u^-)=\Top$; and it satisfies $\nu_+g\le x^+$ iff
$\nu_+g\le\sigma(x^-)^\partial$, iff $\sigma(x^-)\le\nu_-g$. The remaining
entries are the same computation. Both entries of every row and both bounds of
every pair belong to $\chi^\pm\cup B^\pm$, so $\iota$ preserves its solutions.
Finally, $(\iota\rho^\pm)(v^\pm)=(\nu_\mp\rho(v))^\partial=\nu_\pm\rho(v)$ by
\cref{lem:bg-normalization}(1).

(3) If $\iota\sigma=\sigma$, then $\sigma(v^-)=\sigma(v^+)^\partial$ for every $v$,
and \cref{lem:bg-signed} gives $\sigma=\rho^\pm$ for
$\rho(v)=\nu_+\sigma(v^+)$, and the preservation of finiteness.
\end{proof}

\begin{proof}[Proof of \cref{lem:tc-median}]
Write $\operatorname{maj}(L_1,L_2,L_3)$ for the set of words that lie in at
least two of the sets $L_j$. It is symmetric and monotone for inclusion, and
$\operatorname{maj}(L,L,L')=L$.

(1) Symmetry is clear. If $t_j\le t'_j$ for $j=1,2,3$, then
$L^0_{t_j}\subseteq L^0_{t'_j}$ and $L^1_{t'_j}\subseteq L^1_{t_j}$ by
\cref{thm:lf-representation}(2), so the same inclusions hold between the
cones of the medians, and $\tcmed(t_1,t_2,t_3)\le\tcmed(t'_1,t'_2,t'_3)$. The
identity $\tcmed(t,t,u)=t$ follows from $\operatorname{maj}(L,L,L')=L$.

(2) By \cref{thm:lf-representation}(3), with $f$ covariant, a word $iv$ lies
in $L^s_{f(a^j_i)_i}$ iff $v\in L^s_{a^j_i}$, and $\eps$ lies in no cone of a
constructor tree. So the majority of the three cones $L^s_{f(a^j_i)_i}$ is the
union over $i$ of the sets
$i\operatorname{maj}(L^s_{a^1_i},L^s_{a^2_i},L^s_{a^3_i})$, which is
$L^s_{f(m_i)_i}$ for $m_i=\tcmed(a^1_i,a^2_i,a^3_i)$; by
\cref{thm:lf-representation}(1) the two trees are equal. For duality,
$L^s_{t^\partial}=L^{1-s}_t$, and $\operatorname{maj}$ is applied to each sign
separately. For assignments of signed variables,
$(\iota\tcmed(\alpha,\beta,\gamma))(v^\pm)=\tcmed(\alpha(v^\mp),\beta(v^\mp),\gamma(v^\mp))^\partial
=\tcmed((\iota\alpha)(v^\pm),(\iota\beta)(v^\pm),(\iota\gamma)(v^\pm))$.

(3) For $n=0$ every tree is finite. Let $n\ge1$, let $t_1,t_2,t_3$ be finite,
and let $h$ bound the lengths of their positions. A word $w$ of length $h+1$
is a position of no $t_j$, so its longest prefix in the domain of $t_j$ is a
leaf, labeled $\Bot$ or $\Top$; thus $w\in L^0_{t_j}\cup L^1_{t_j}$ for each
$j$. Two of the three trees put $w$ in the same cone $L^s$, so $w$ lies in the
cone $L^s$ of the median, and some prefix of $w$ is a leaf of the median.
Every proper prefix of a position is labeled $f$, so no word of length $h+2$
is a position of the median, which is therefore finite.

(4) Let $\alpha,\beta,\gamma$ solve a covariant system together with covariant
ground bounds. A literal $u\le v$ holds in their median $\tcmed(\alpha,\beta,\gamma)$
by monotonicity. For $f(a_i)_i\le u$, monotonicity gives
$\tcmed(f(\alpha(a_i))_i,f(\beta(a_i))_i,f(\gamma(a_i))_i)\le
\tcmed(\alpha,\beta,\gamma)(u)$, and the left side is
$f(\tcmed(\alpha,\beta,\gamma)(a_i))_i$ by (2); the upper constructor literals are
symmetric. The extreme literals hold since $\tcmed(\Bot,\Bot,\Bot)=\Bot$ and
$\tcmed(\Top,\Top,\Top)=\Top$. For a lower bound $G\le p$,
$G=\tcmed(G,G,G)\le\tcmed(\alpha,\beta,\gamma)(p)$, and upper bounds are
symmetric.
\end{proof}

\subsection{Covariant Finite-Bound Pinning}\label{app:tc-covariant}

In this subsection $f$ is covariant and $n\ge1$; for $n=0$,
\cref{lem:tc-pinning-covariant} is trivial because every tree is finite. We use
the closure $D^{\mathrm{cov}}_\chi$, the path bounds $\Lo$ and $\Up$, label clashes and cycle
clashes of a covariant system without copy literals, as defined in
\cref{app:bg-sat}, and \cref{thm:bg-nsse-sat}(1): such a system has a
solution iff it has no label clash, and a finite solution iff it has neither a
label clash nor a cycle clash.

\begin{proof}[Proof of \cref{lem:tc-pinning-covariant}]
Let $F$ be a covariant system over $V$ with a finite solution $\delta$, let $B$
be a finite set of ground bounds on $V$, and let $\rho$ be a solution of
$F\cup B$.

\emph{Copy literals.} Let $F_0=F^\flat$ be obtained by
\cref{lem:bg-sandwich}, applied to every copy literal of $F$. The variables of
$F_0$ are called \emph{old}. By \cref{lem:bg-sandwich}, $\delta$ extends to a
finite solution of $F_0$, $\rho$ extends to a solution of $F_0\cup B$, and
every solution of $F_0\cup B$ restricts to a solution of $F\cup B$.

\emph{Encoding the bounds.} For a finite ground tree $g$ and a variable $x$
define systems $L(g,x)$, expressing $g\le x$, and $U(x,g)$, expressing $x\le
g$, with fresh variables:
\[
\begin{aligned}
 L(\Bot,x)&=\varnothing,\quad L(\Top,x)=\{\Top\le x\},\quad
 L(f(g_i)_i,x)=\{f(y_i)_i\le x\}\cup\textstyle\bigcup_iL(g_i,y_i),\\
 U(x,\Top)&=\varnothing,\quad U(x,\Bot)=\{x=\Bot\},\quad
 U(x,f(g_i)_i)=\{x\le f(y_i)_i\}\cup\textstyle\bigcup_iU(y_i,g_i).
\end{aligned}
\]
Each fresh variable stands for a nonempty position $\pi$ of $g$, and its
\emph{rank} is $\len\pi$; old variables have rank $0$. By induction on $g$, an
assignment that satisfies $L(g,x)$ satisfies $g\le x$, since $f(g_i)_i\le
f(y_i)_i\le x$ by monotonicity of the covariant constructor; conversely, an
assignment that satisfies $g\le x$ extends to $L(g,x)$ by giving each fresh
variable the subtree of $g$ at its position. The same holds for $U(x,g)$. Let
$F''$ be $F_0$ together with the encodings of all bounds of $B$, with distinct
fresh variables. Then $F''$ has a solution, which extends one of $F_0\cup B$,
and every solution of $F''$ restricts to a solution of $F_0\cup B$. The fresh
variables of the systems $L$ are \emph{sources}, those of the systems $U$ are
\emph{sinks}, and the literals of $F''$ outside $F_0$ are \emph{new}. The
system $F''$ has no copy literals, and by construction:
\begin{enumerate}[label=(O\arabic*),leftmargin=*]
\item the literals of $F_0$ contain only old variables;
\item the new lower literals are $f(y_i)_i\le x$ and $\Top\le x$, where $x$ is
  old or a source and every $y_i$ is a source of rank one more than $x$; the
  new upper literals are $x\le f(y_i)_i$ and $x=\Bot$, where $x$ is old or a
  sink and every $y_i$ is a sink of rank one more than $x$.
\end{enumerate}
In particular, the children of a lower constructor literal are all old or all
sources, and the children of an upper constructor literal are all old or all
sinks.

\emph{Separation.} Write $D=D^{\mathrm{cov}}_{F''}$, $\Lo$, $\Up$ for the
closure and the path bounds of $F''$, and $D_0=D^{\mathrm{cov}}_{F_0}$, $\Lo_0$,
$\Up_0$ for those of $F_0$. We show by
induction on the derivation of $D(a,b)$:
(a) if $b$ is a source, then $a=b$;
(b) if $a$ is a sink, then $b=a$;
(c) if $a$ and $b$ are old, then $D_0(a,b)$.
Reflexivity is clear. Let $D(a,b)$ follow from $D(a,m)$ and $D(m,b)$. For (a),
the induction hypothesis gives $m=b$ and then $a=m$; (b) is symmetric. For (c),
$m$ is not a source, since (a) would give $a=m$, and not a sink, since (b)
would give $b=m$; so $m$ is old, and $D_0(a,m)$ and $D_0(m,b)$. Let
$D(a_i,b_i)$ follow by \cref{eq:bg-decomp} from literals $f(a_j)_j\le u$ and
$v\le f(b_j)_j$ and $D(u,v)$. Then $b_i$ is old or a sink and $a_i$ is old or
a source, so (a) and (b) hold vacuously. If $a_i$ and $b_i$ are old, the two
literals have old children; by (O2) they are not new, so they are literals of
$F_0$, and $u,v$ are old by (O1). By the induction hypothesis $D_0(u,v)$, hence
$D_0(a_i,b_i)$. Conversely $D_0\subseteq D$, so $D$ and $D_0$ agree on old
variables.

\emph{Paths.} We show by induction on $\eta$:
(d) if $\Lo(\eta,a,s)$ and $s$ is a source, then $a$ is a source of rank
$\operatorname{rank}(s)+\len\eta$;
(e) if $\Up(\eta,t,b)$ and $t$ is a sink, then $b$ is a sink of rank
$\operatorname{rank}(t)+\len\eta$;
(f) if $a,b$ are old, then $\Lo(\eta,a,b)$ implies $\Lo_0(\eta,a,b)$ and
$\Up(\eta,a,b)$ implies $\Up_0(\eta,a,b)$.
For $\eta=\eps$ use (a)--(c). For (d), by \cref{eq:bg-paths},
$\Lo(i\eta,a,s)$ gives a literal $f(a'_j)_j\le z$ with $D(z,s)$ and
$\Lo(\eta,a,a'_i)$. By (a), $z=s$, so the literal has a source as its root; by
(O1) it is new, and by (O2) $a'_i$ is a source of rank
$\operatorname{rank}(s)+1$. The induction hypothesis applies to
$\Lo(\eta,a,a'_i)$. (e) is symmetric, using (b). For (f), $\Lo(i\eta,a,b)$
gives $f(a'_j)_j\le z$ with $D(z,b)$ and $\Lo(\eta,a,a'_i)$. If the literal
were new, $a'_i$ would be a source, and (d) would make $a$ a source. So it is a
literal of $F_0$, $z$ and $a'_i$ are old, $D_0(z,b)$ by (c), and
$\Lo_0(\eta,a,a'_i)$ by the induction hypothesis. The case of $\Up$ is
symmetric, using (e).

\emph{No cycle clash.} Suppose that $F''$ had a cycle clash
\cref{eq:bg-cycle-clash}: $\eta\ne\eps$, $\Lo(\eta,a,a)$, $D(a,b)$ and
$\Up(\eta,b,b)$. If $a$ were a source, (d) would give
$\operatorname{rank}(a)=\operatorname{rank}(a)+\len\eta$. If $b$ were a sink,
(e) would give the same contradiction for $b$. If $a$ were a sink, (b) would
give $b=a$, a sink; if $b$ were a source, (a) would give $a=b$, a source. So
$a$ and $b$ are old, and by (c) and (f) the clash is a cycle clash of $F_0$.
But $F_0$ has a finite solution, so it has no cycle clash.

\emph{Conclusion.} The system $F''$ has a solution, so it has no label clash,
and it has no cycle clash. Hence it has a finite solution, whose restriction is
a finite solution of $F_0\cup B$ and then of $F\cup B$.
\end{proof}

The proof extends the argument of
\nssecite{Theorem ``Finite countermodel transfer''}, where the bounds are the
pin paths of a single word, to arbitrary finite ground bounds. The finite solution
of $F$ is used only to exclude cycle clashes among old variables; label clashes
are excluded by the solution of $F\cup B$.

\subsection{Closure under Intersection}

\begin{proof}[Proof of \cref{prop:ch-meet}]
(1) Let $\mathcal M$ be a nonempty family of models of $\chi$ and
$M=\bigcap\mathcal M$. If the premise of a rule lies in $M$, it lies in every
member of $\mathcal M$, hence so does its conclusion; so $M$ is closed under
the rules. It contains $\Fact_\chi$, since every member does, and it is
consistent, since it is contained in a member. For the projected family, let
$S_j=U_j|_X$ with $U_j\in\Mod(\psi)$ for $j$ in a nonempty index set; then
$\bigcap_jU_j$ is a model of $\psi$ with visible part $\bigcap_jS_j$. For
assignments, $S_{\rho\leafmeet\tau}=S_\rho\cap S_\tau$ by definition, so
\cref{lem:lf-clauses} turns the first two claims into the closure of
solutions and of extensible visible assignments under $\leafmeet$: if $\rho$
and $\tau$ extend to solutions $\rho'$ and $\tau'$ of $\psi$, then
$\rho'\leafmeet\tau'$ is a solution of $\psi$ that extends
$\rho\leafmeet\tau$.

(2) Commutativity, associativity and idempotence hold coordinatewise. If
$t\le_c t'$, then $L^0_t\cap L^0_u\subseteq L^0_{t'}\cap L^0_u$ and
$L^1_{t'}\cap L^1_u\subseteq L^1_t\cap L^1_u$, so
$t\leafmeet u\le_c t'\leafmeet u$ by \cref{thm:lf-representation}(2). For the
constructor, \cref{thm:lf-representation}(3) gives
$L^s_{f(t_i)_i}=\bigcup_iiL^{s\oplus c(i)}_{t_i}$, a union of sets of words
with distinct first letters. Intersection therefore distributes over these
unions, and the leaf pairs of both sides of the identity agree; by
\cref{thm:lf-representation}(1) the trees agree.

For the median, observe that the pairs
$(L^0_t\cup L^0_u,\ L^1_t\cap L^1_u)$ and
$(L^0_t\cap L^0_u,\ L^1_t\cup L^1_u)$ are disjoint and extension closed: for
instance, a word of $L^0_t$ in $L^1_t\cap L^1_u$ would lie in both cones of
$t$. By \cref{thm:lf-representation}(1,2) they are the leaf pairs of the least
upper bound $t\vee u$ and the greatest lower bound $t\wedge u$ of $t$ and $u$
for $\le_c$. So the trees form a lattice whose joins and meets act on the
cone of sign $0$ by union and intersection and on the cone of sign $1$ by
intersection and union; in particular this lattice is distributive. The
median $(t\wedge u)\vee(u\wedge v)\vee(v\wedge t)$ of three trees therefore
has, as each of its two cones, the set of words that lie in the corresponding
cones of at least two of $t$, $u$ and $v$. For $v=\Omega$, whose cones are
empty, these are the words that lie in the cones of both $t$ and $u$, so the
median is $t\leafmeet u$.

(3) Regular languages are closed under intersection, so the first claim
follows from \cref{thm:lf-representation}(4). The cones of $\Bot\leafmeet\Top$
are $\varnothing\cap\Sigma^*$ and $\Sigma^*\cap\varnothing$, so it is $\Omega$,
all of whose words are positions labeled $f$; for $n\ge1$ there are infinitely
many.
\end{proof}
